\documentclass[10pt,a4paper]{amsart}
\usepackage[margin=19mm]{geometry}
\usepackage{amsmath,amssymb,mathtools}
\usepackage[scr=rsfs,cal=euler]{mathalfa}
\usepackage{xfrac}
\usepackage[unicode,hidelinks,bookmarks=false]{hyperref}
\newcommand{\ie}{i.e.\ }
\newcommand{\cf}{cf.\ }
\newcommand{\resp}{respectively\ }
\newcommand{\Subsection}{\subsection}
\providecommand{\nobreakdash}{\nobreak\hbox{-}}
\setlength{\emergencystretch}{2em}
\multlinegap0pt
\usepackage{amssymb}
\newtheorem{theorem}{Theorem}
\newtheorem{proposition}{Proposition}
\newcommand{\frc}[1]{\{ #1 \}}
\title{Beatty solutions of almost Golomb functional equations}
\author{Benoit Cloitre}
\date{Source: arXiv:2604.10822v4 (CC BY 4.0)}
\begin{document}
\maketitle

\noindent\emph{Source and licence.} Beatty solutions of almost Golomb functional equations. Version arXiv:2604.10822v4 (CC BY 4.0), distributed by arXiv under the Creative Commons Attribution 4.0 International licence, http://creativecommons.org/licenses/by/4.0/, as recorded in the arXiv licence metadata for this exact version and shown on the abstract page https://arxiv.org/abs/2604.10822v4. Full licence notice: \texttt{licenses/arxiv-2604.10822-CC-BY-4.0.txt}.\par\smallskip

\noindent\textit{Editorial scope.} Counted unit: the article's proposition prop:defect\_substitution (source0.tex lines 2693--2703). The article's complete original proof of it is delivered with it (source0.tex lines 2705--2792, one \texttt{proof} environment of 88 lines). No mathematics is written, repaired or re-derived: the statement and the proof are the article's own lines, in the article's own notation, and no retained line is substituted or re-flowed.  Retained uncounted as the source-local context the proof needs: lines 2573--2591.  Nothing was omitted from any retained span, so the package carries no omission record.  No retained line contains a citation, so the wrapper carries no bibliography and no bibliography entry.  No retained line contains a figure environment, an image-inclusion command or drawing code, so no image is delivered and the package declares no asset dependency.  Editorial formatting only, in addition: an AMS \texttt{amsart} preamble that restates the article's own declarations and macros used by the retained lines, namely the environments \texttt{theorem}, \texttt{proposition} on separate counters, together with the standard packages those lines need.  The retained environments print in this document as Theorem 1, Proposition 1 in order of appearance, whereas the article prints the same items with the numbering of its own full document.  The complete native source of the article as distributed in the arXiv e-print for this exact version is bundled unchanged as \texttt{sources/198295/source0.tex}.
\begin{theorem}\label{thm:absorption}
For $d=2(\sqrt{2}{-}1)$ and $n\ge 2$:
\[
  \mu(n)=1\iff
  \theta_{n-1}\in[1-c,\,1/2].
\]
Consequently:
\begin{enumerate}
\item[\textup{(a)}] Every $n\in\mathcal{D}$ is a repeat
  of~$a_{\mathrm{R}}$, i.e.
  $a_{\mathrm{R}}(n{+}1)=a_{\mathrm{R}}(n)$.
\item[\textup{(b)}] The set~$\mathcal{D}$ has natural
  density
\[
  \rho(\mathcal{D})=c-\tfrac{1}{2}
  =\frac{\sqrt{2}-1}{2}\approx 0.2071.
\]
\end{enumerate}
\end{theorem}

\begin{proposition}\label{prop:defect_substitution}
The gap sequence of~$\mathcal{D}$ on the alphabet
$\{3,4,7\}$ has letter frequencies
\[
  f_{3}=\sqrt{2}-1,\quad f_{4}=3-2\sqrt{2},\quad f_{7}=\sqrt{2}-1,
\]
matching Theorem~\ref{thm:absorption}(b), and the
average gap is $2(\sqrt{2}+1)$. The gap sequence
of~$\mathcal{D}$ is
\href{https://oeis.org/A395253}{\texttt{A395253}}.
\end{proposition}

\begin{proof}
The defect set is
$\mathcal{D}=\{n\ge 2:\theta_{n-1}\in J\}$ with
$J=[1-c,\,1/2]$, so $|J|=c-1/2=(\sqrt{2}-1)/2$.
For each return time $i\ge 1$, let
\[
  K_{i}=\{\theta\in J: R_{c}^{i}\theta\in J
  \text{ and } R_{c}^{k}\theta\notin J
  \text{ for }1\le k<i\},
\]
where $R_{c}\colon\theta\mapsto\frc{\theta+c}$.
We compute the sub-intervals directly from the
iterates. For $\theta\in J=[1-c,1/2]$, a direct reduction
modulo~$1$ gives
\[
  R_{c}\theta=\theta+c-1,\qquad
  R_{c}^{2}\theta=\theta+2c-1,\qquad
  R_{c}^{3}\theta=\theta+3c-2
\]
for $\theta\in J$. Neither $R_{c}\theta\in[0,c-1/2]$
nor $R_{c}^{2}\theta\in[c,2c-1/2]$ meets~$J$, so
every return time is at least~$3$.

Return time~$3$ occurs when
$\theta+3c-2\in[1-c,1/2]$, equivalently
$\theta\in[3-4c,\,5/2-3c]$. Since $1-c>3-4c$
and $5/2-3c<1/2$, intersection with~$J$ gives
\[
  K_{3}=[1-c,\,5/2-3c],\qquad
  |K_{3}|=(5/2-3c)-(1-c)=3/2-2c=\tfrac{3-2\sqrt{2}}{2}.
\]

For $\theta\in J\setminus K_{3}=(5/2-3c,\,1/2]$, the
return time exceeds~$3$. The next iterate
$R_{c}^{4}\theta=\theta+4c-3$ lies in~$J$ when
$\theta\in[4-5c,\,7/2-4c]$. Intersection with
$(5/2-3c,\,1/2]$ gives
\[
  K_{4}=[4-5c,\,1/2],\qquad
  |K_{4}|=1/2-(4-5c)=5c-7/2=\tfrac{5\sqrt{2}-7}{2}.
\]

The remainder
$K_{7}=(5/2-3c,\,4-5c)=J\setminus(K_{3}\sqcup K_{4})$
has length
\[
  |K_{7}|=(4-5c)-(5/2-3c)=3/2-2c=\tfrac{3-2\sqrt{2}}{2},
\]
For $\theta\in K_{7}$, the fifth and sixth iterates satisfy
\[
  R_{c}^{5}\theta=\theta+5c-3\in\bigl(2c-\tfrac12,\,1\bigr),
  \qquad
  R_{c}^{6}\theta=\theta+6c-4\in\bigl(3c-\tfrac32,\,c\bigr),
\]
both disjoint from $J=[1-c,1/2]$, since $3c-\tfrac32>\tfrac12$.
On the other hand,
\[
  R_{c}^{7}\theta=\theta+7c-5\in\bigl(4c-\tfrac52,\,2c-1\bigr)
  \subset J,
\]
so the first-return time is $7$ throughout $K_{7}$, and
$J=K_{3}\sqcup K_{7}\sqcup K_{4}$ is the exact first-return
partition.

The orbit $(\theta_{n-1})_{n\ge 2}$ equidistributes
in $[0,1)$ by Weyl's theorem, so each $K_{i}$ is
visited with density~$|K_{i}|$, and the frequency
of gap value~$i$ among defect positions is
\[
  f_{i}=\frac{|K_{i}|}{|J|}.
\]
Substituting the lengths and using
$(\sqrt{2}-1)(\sqrt{2}+1)=1$ to rationalize yields
\[
  f_{3}=f_{7}=\sqrt{2}-1,\qquad
  f_{4}=3-2\sqrt{2}.
\]
The average gap is
\[
  3f_{3}+4f_{4}+7f_{7}
  =3(\sqrt{2}{-}1)+4(3{-}2\sqrt{2})+7(\sqrt{2}{-}1)
  =2\sqrt{2}+2
  =2(\sqrt{2}{+}1),
\]
which equals the reciprocal of the density
$\rho(\mathcal{D})=(\sqrt{2}-1)/2$ established in
Theorem~\ref{thm:absorption}(b).
\end{proof}

\end{document}
